\begin{papercard}{polson1992}
\cf{Stages} \stagetag{W}. The theory paper and detailed procedure of the \emph{second} cognitive walkthrough version (CW2). It covers the exploratory-learning model, the three-section per-action form with percentage estimates, a fully worked phone example, and early validation. \emph{Note:} the appended copy is the ICS Technical Report 91-01 (an OCR'd scan), the preprint of the 1992 IJMMS article, so page numbers refer to the report.

\cf{Rationale} A hand simulation of the user's cognitive activities, adapted from software-engineering design walkthroughs, to check that the user can easily learn the intended tasks \pc{polson1992}{3}{The method we present involves hand simulation of the cognitive activities of a user}. Focus: ease of learning by exploration, before an implementation or mock-up exists \pc{polson1992}{4}{The method gives designers the}. \textbf{Inputs}: a detailed interface description, a task scenario, explicit assumptions about users and context, and a correct action sequence \pc{polson1992}{5}{The input to a cognitive walkthrough session includes a detailed design description}. \textbf{Evidence rule}: claims that a step is fine \emph{or} problematic must be supported by theory, data or experience \pc{polson1992}{5}{a given step would not cause any difficulties must be supported by}.

\cf{Theory} Exploration is ``guessing what to do using cues provided by the system'' \pc{polson1992}{6}{guess what to do using cues provided by the system}. Built on Norman's action cycle and Kintsch's construction--integration model. Goals are propositions linked to knowledge, objects and actions, and activation spreads to actions \pc{polson1992}{7}{We have developed our model within the framework of Kintsch}. \emph{Label following}: an action is chosen because its label shares terms with an active goal \pc{polson1992}{8}{Here an action, such as pressing a button, is chosen because there is a label}. The goal--action path can break: no shared terms, an unclear label placement, missing knowledge, or competing labels \pc{polson1992}{9}{there could be two or more labels that seem to be}. The procedure makes knowledge assumptions explicit (e.g.\ ``1/0'' vs.\ ``on/off'') \pc{polson1992}{9}{The walkthrough procedure helps to make these assumptions explicit}. Goals are built from interface prompts and feedback \pc{polson1992}{10}{The goal generation process will be driven almost exclusively by prompts}. ``And-then'' goal structures explain \emph{supergoal kill-off}: a final terminator such as ENTER is forgotten when a similar subgoal is done \pc{polson1992}{13}{We call this the supergoal kill-off phenomenon}.

\cf{Phase 1: preparation (start-up sheet)} Two phases, preparation and evaluation, with a printed form that works ``like an income tax form'' for non-psychologists \pc{polson1992}{14}{The approach is similar to}. \emph{Tasks}: task selection is outside the method's scope, but tasks should be what users actually want to do and should include sequences of basic tasks, because transitions cause problems \pc{polson1992}{16}{How tasks should be selected is outside the scope of the method}. \emph{Task description}: from the first-time user's view, in realistic, non-system terms, with the starting state noted \pc{polson1992}{16}{It is important to describe the tasks in}. \emph{Action sequence}: the \emph{best} sequence the interface offers, since problems on an inferior path teach little. All acceptable alternatives should be analysed \pc{polson1992}{18}{the sequence be the best offered by the interface}. \emph{Granularity}: start finely divided and collapse only when no problem will be hidden \pc{polson1992}{18}{start by thinking of actions being very finely divided}. \emph{Users}: a relatively uniform population with its background knowledge \pc{polson1992}{19}{a relatively uniform population of potential users needs to be identified}. \emph{Initial goals}: the goals users are \emph{likely} to have, with their source (task, screen, knowledge) and percentages for alternative goal structures \pc{polson1992}{19}{Each goal should be annotated as to whether it is based on the task description}.

\cf{Phase 2: per-action form} Three sections per action, with most items estimating the percentage of users who fail (0/25/50/75/100) to rank severity \pc{polson1992}{20}{Most items also ask the analyst to estimate the}. \emph{Section 1 (goals)}: the correct goals and the mismatch with likely goals. Problems are split into \emph{goal problems} (trying the wrong thing) and \emph{action problems} (right goal, cannot find how) \pc{polson1992}{20}{In a goal problem, the user is trying to do the wrong thing}. How the correct goals are specified shifts problems between the two types, reflecting design philosophy \pc{polson1992}{22}{Generally, a designer can work with more}. \emph{Section 2 (action)}: availability, label, link of label to action, link of label to goal, no label, wrong choices, time-out, hard to do (items 2.1--2.8) \pc{polson1992}{23}{it obvious that the correct action is a possible choice here}. \emph{Section 3 (goal modification)}: quit or back up, accomplished goals, incomplete goals that look accomplished, and-then check, and new goals from prompts or other sources (items 3.1--3.6) \pc{polson1992}{25}{Incomplete goals that look accomplished}. Section 3 percentages are inputs to the next step's goal analysis, not problems in themselves \pc{polson1992}{27}{The percentages of users forming or not forming certain goals should not be taken}.

\cf{Worked example} ``Forward all my calls'', when calls are already forwarded elsewhere (seven actions). The analysts assign a recorder, estimate that 75\,\% of users lack the ``clear forwarding'' goal, find that the CLEAR TRANSFER label competes with the correct one (25\,\%), and note uninformative beep feedback \pc{polson1992}{27}{One of the analysts is assigned to be the recorder}, \pc{polson1992}{32}{and the analysts fear that CLEAR TRANSFER ALL CALLS is as good a match}. They resolve a debated fact \emph{by experiment} \pc{polson1992}{32}{By experiment they determine that forwarding really is cancelled}. The goal structure is revised retrospectively when a later step reveals a hidden subgoal \pc{polson1992}{33}{they decide they will go back and add a new subgoal}. Three remedies: training, a prompt, or removing the need, the last being best \pc{polson1992}{36}{providing training so that users will bring the right goal structure to the task}. After a fix, the \emph{whole} walkthrough must be reviewed \pc{polson1992}{37}{since a changed label on the template might incorrectly match a goal at any point in the action sequence}.

\cf{Evidence} CW1 found about half of observed errors with a false-alarm rate of almost 75\,\% \pc{polson1992}{37}{The false alarm rate, the proportion of errors identified in the walkthrough}. Student projects: 50\,\% of think-aloud errors, about 30\,\% overall and over 70\,\% of the most serious ones, and all problems without false alarms on one simple task \pc{polson1992}{37}{Testing with the voicemail directory showed that 50}. Jeffries et al.: HE 105, usability testing 31, guidelines 35, CW 35. CW was usable by developers but missed ``general and recurring problems'' \pc{polson1992}{39}{Disadvantages of the walkthrough technique included the need for a methodology to define user tasks}.

\cf{Variability between analysts} Sources: poor procedural control (different paths, forgotten steps), different care in specifying goals, uncertain facts not derivable from theory, and differences in noticing \pc{polson1992}{40}{Some differences are traceable to poor}. Earlier evaluations were done individually to study interrater reliability. Group analysis was ``a very different experience'' that integrates viewpoints \pc{polson1992}{40}{We have done a few analyses as a group and find it a very different experience}. The method is tedious \pc{polson1992}{40}{Another conclusion from all experiences with the method is that it is tedious}.

\cf{Positioning} CW identifies problems only in task context and links mental processes to specific prompts and cues \pc{polson1992}{48}{The cognitive walkthrough identifies problems only in}. It does not cover enjoyment or general guidelines \pc{polson1992}{46}{The cognitive walkthrough examines only whether users will be}. User testing finds problems CW misses, so CW precedes and focuses testing and helps diagnose goal vs.\ execution problems \pc{polson1992}{47}{user testing will find problems that the walkthrough misses}. Other uses: design rationale, self-critique, and non-walk-up systems where background knowledge matters and menus are searched \pc{polson1992}{50}{the intended path through an interface may include searching for information}.

\cf{Use for this thesis} The most detailed form for coding a privacy walkthrough step. Items 2.1--2.6 (availability, label, label--action, label--goal, wrong choices) map directly onto finding a privacy setting or policy section, and item 3.3 (``looks accomplished but is not'') fits consent and deletion flows that suggest completion without effect. Three practices to adopt: (1) the evidence rule, where every pass/fail claim cites a screenshot or a stated assumption; (2) the goal-problem vs.\ action-problem distinction as a codebook dimension; (3) the reported \emph{sources of analyst variability} as the rationale for an independent second coder and a fixed path. Percentage estimates are optional and should be framed as judgements, not data.
\end{papercard}
